\section{多头不是冗余而是分工}
\label{sec:14-Deep-Learning/05-Attention-and-Sequence-Models/02}

你在做一组消融实验：把一个十二层模型每层的八个头合并成一个头，总宽度仍然是 \(d=512\)，只是这唯一的头把 \(d_k\) 从 64 提到 512。参数量一个不多一个不少，浮点运算量也几乎相同，数据、步数、学习率全部照旧。验证困惑度掉了一截，换三个随机种子都一样。

同样的参数、同样的算力，效果却不同，说明差别不在容量而在结构。把单头输出的形状写出来，原因就摆在式子里。

\subsection{一个头只能给出一种对齐}

\hyperref[sec:14-Deep-Learning/05-Attention-and-Sequence-Models/01]{``Self-Attention 是内容决定的加权读取''}给出的输出是

\[
z_t = \sum_{s=1}^{n} A_{ts} v_s,\qquad A_{ts}\ge 0,\quad \sum_s A_{ts}=1.
\]

权重非负且行和为一，所以 \(z_t\) 落在 Value 行向量 \(v_1,\ldots,v_n\) 的凸包里。这不是一句技术细节：它意味着单个头对每个查询位置只能交出\emph{一个}混合比例，而混合是不可逆的——一旦把两份信息按 0.5 与 0.5 加起来，下游拿到的就是它们的平均，再也拆不开。

考虑一个真实需求。位置 \(t\) 上的词同时需要两样东西：一是紧邻的前一个词（决定局部句法角色），二是句首那个远处的名词（决定代词指向谁）。单头有且只有一行权重可以分配。它把权重压在邻居上，指代信息丢失；分给远处名词，句法信息被稀释；对半开，则 \(z_t\) 是两个向量的中点，下游的前馈层看到的是一个既不像句法特征也不像指代特征的混合物。除非 \(v_s\) 恰好分布在互相正交的子空间里，否则这两份信息在加权求和的那一步就被压扁了——而单头没有任何机制保证这种正交性。

多头做的恰恰是把这种正交性写进结构。第 \(i\) 个头用自己的 \(W_Q^{(i)},W_K^{(i)}\in\R^{d\times d_h}\) 与 \(W_V^{(i)}\in\R^{d\times d_h}\)（\(d_h=d/h\)）算出 \(Z^{(i)}\in\R^{n\times d_h}\)，然后

\[
Z = \bigl[\,Z^{(1)}\;Z^{(2)}\;\cdots\;Z^{(h)}\,\bigr]\,W_O,
\qquad W_O\in\R^{d\times d}.
\]

拼接是一次硬性的坐标分块：第 \(i\) 个头的读取结果只写进第 \(i\) 段坐标，任何两个头的输出天然落在互不重叠的子空间里。指代那份信息和句法那份信息各占一段，先并排放好，再由 \(W_O\) 决定怎么混。混合发生在 \(W_O\) 里，是可学习的、对下游任务负责的；单头的混合发生在 softmax 加权那一步，是被迫的、由匹配分数顺带决定的。这是两种完全不同的信息流。

\subsection{参数量为什么原地不动}

八个头看上去是八套投影，很容易误以为参数翻了八倍。数一遍就清楚：每个头的 \(W_Q^{(i)}\) 是 \(d\times d_h\)，\(h\) 个头拼起来正好是 \(d\times (h\cdot d_h)=d\times d\)。\(K\) 与 \(V\) 同理。加上 \(W_O\)，总量恒为 \(4d^2\)，与 \(h\) 无关。

\[
h \times \underbrace{3\,d\,\frac{d}{h}}_{\text{单头三组投影}} + \underbrace{d^2}_{W_O} \;=\; 4d^2 .
\]

浮点运算量也基本持平：\(h\) 个 \(n\times n\) 分数矩阵，每个的内积维度是 \(d/h\)，乘起来仍是 \(O(n^2 d)\)。所以多头不是用更多参数换更强表达，而是在同一笔预算里换了一种切分方式。省下来的东西没有，改变的只是这 \(4d^2\) 个数被组织成几个独立的匹配通道。

\subsection{每个头的匹配是一个低秩双线性型}

分工的能力不是白来的，限制藏在秩里。把分数矩阵完全展开：

\[
S^{(i)} = \frac{1}{\sqrt{d_h}}\,X W_Q^{(i)} \bigl(X W_K^{(i)}\bigr)^{\trans}
        = \frac{1}{\sqrt{d_h}}\,X M^{(i)} X^{\trans},
\qquad M^{(i)} \eqdef W_Q^{(i)} W_K^{(i)\trans} \in \R^{d\times d}.
\]

一个头的全部匹配行为压缩在 \(M^{(i)}\) 这一个矩阵里，它是 \(d\) 维表示空间上的一个双线性型：\(x_t\) 与 \(x_s\) 的匹配分数就是 \(x_t^{\trans}M^{(i)}x_s\)。

\begin{proposition}
\(\rank M^{(i)}\le d_h = d/h\)，从而 \(\rank S^{(i)}\le d/h\)。特别地，当 \(d/h < n\) 时，存在无论 \(W_Q^{(i)},W_K^{(i)}\) 怎么取都无法实现的 \(n\times n\) 分数矩阵。
\end{proposition}

证明骨架：\(M^{(i)}\) 是 \(d\times d_h\) 与 \(d_h\times d\) 两个矩阵的乘积，秩不超过公共维度 \(d_h\)，这是矩阵乘法的秩不等式，不需要任何附加条件。\(S^{(i)}\) 是 \(X M^{(i)} X^{\trans}\)，左右各乘一个矩阵只会让秩不增。条件对账：结论用到的唯一事实是投影的中间维度等于 \(d_h\)；一旦允许 \(W_Q,W_K\) 的列数大于 \(d_h\)，上界随之放宽——这正是单头 \(d_k=512\) 相对八头 \(d_h=64\) 的全部优势所在。注意上界针对的是 softmax \emph{之前}的分数矩阵；softmax 逐元素作用，会破坏低秩结构，所以 \(A^{(i)}\) 本身的秩可以更高，但它能表达的模式仍被 \(S^{(i)}\) 的低秩几何所约束。

这个秩上界把头数变成一个真正的权衡量。\(h\) 增大，独立的对齐通道变多，但每条通道的匹配是 \(d/h\) 维空间里的一个低秩形式，能区分的方向越来越少；\(h\) 减小，每个头的匹配精细，但可同时表达的对齐关系变少。\(d=512\)、\(h=64\) 时每个头只剩 8 维——在 8 维里做 Query-Key 点积，能编码的匹配规则相当粗糙。这与 \hyperref[sec:03-Linear-Algebra/07-SVD-and-Data-Applications/03]{``低秩近似：用少数方向保留主要结构''}讨论的是同一件事：低秩不是缺陷，是一种被强加的归纳偏置，它决定了哪些结构可以被保留、哪些必然被抹掉。

\begin{center}
\begin{tikzpicture}[scale=0.42]
% ---- 头 A：邻接带状 ----
\begin{scope}
\foreach \r in {1,...,8}{
  \foreach \c in {1,...,8}{
    \pgfmathtruncatemacro{\d}{abs(\r-\c)}
    \pgfmathtruncatemacro{\v}{max(3,72-30*\d)}
    \fill[black!\v] ({\c-1},{-\r}) rectangle ({\c},{1-\r});
  }
}
\draw[gray!50] (0,0) grid (8,-8);
\draw[black!70] (0,0) rectangle (8,-8);
\node at (4,1.0) {\footnotesize 头 A：带状（邻接）};
\node[rotate=90] at (-1.1,-4) {\footnotesize 查询 \(t\)};
\node at (4,-9.1) {\footnotesize 键 \(s\)};
\end{scope}
% ---- 头 B：远距离稀疏对应 ----
\begin{scope}[shift={(11,0)}]
\foreach \r in {1,...,8}{
  \foreach \c in {1,...,8}{
    \fill[black!4] ({\c-1},{-\r}) rectangle ({\c},{1-\r});
  }
}
\foreach \r/\c/\v in {1/1/40, 2/7/68, 3/8/60, 4/2/62, 5/1/55, 6/8/70, 7/3/64, 8/5/66,
                      2/8/22, 4/1/20, 6/2/18, 8/1/24}{
  \fill[black!\v] ({\c-1},{-\r}) rectangle ({\c},{1-\r});
}
\draw[gray!50] (0,0) grid (8,-8);
\draw[black!70] (0,0) rectangle (8,-8);
\node at (4,1.0) {\footnotesize 头 B：稀疏（远距离）};
\node at (4,-9.1) {\footnotesize 键 \(s\)};
\end{scope}
% ---- 共同输入 ----
\draw[->,>=stealth,black!60] (8.4,-4) -- (10.6,-4);
\node at (9.5,-3.4) {\footnotesize 同一份 \(X\)};
\end{tikzpicture}
\end{center}

\emph{同一个句子、同一层，两个头给出的注意力权重矩阵。左边几乎全部质量压在对角带上，右边把质量押在少数远处位置。输入一模一样，读取规则完全不同——这就是所谓分工：它不是模型被告知要做的事，而是不同子空间里独立匹配的结果。}

图里两块热度图的稀疏程度差别很大，但它们受同一个秩上界约束。带状那种模式对低秩其实相当友好：只依赖 \(t-s\) 的分数矩阵是一个 Toeplitz 结构，用少数几个方向就能近似得不错。稀疏的远距离对应则难得多——要让第 2 行的质量准确落在第 7 列而不溅到别处，需要 \(M^{(i)}\) 在很多方向上都能分辨，8 维空间往往不够。头维度砍得太狠时，先失去的通常是右边那种模式。

\subsection{分工是观察到的，不是证明出来的}

到这里为止的论证只说明了一件事：多头\emph{允许}不同的对齐关系并存。它没有说每个头一定会各占一个语义位置。

大量探测实验确实报告过可命名的头——有的头几乎总是关注前一个词，有的稳定地把代词连到先行词，有的对应句法依存边。这些是在特定模型、特定数据上观察到的经验规律，不是定理。同一架构换个随机种子重训，头的编号与功能的对应关系会整体洗牌；同一个头在不同输入分布下的行为也未必一致。

更需要坦白的是冗余。剪枝实验反复显示，训练好的模型里相当比例的头可以直接移除而下游指标几乎不动，某些层甚至只留一个头就够。这不奇怪：\(h\) 个头是并行的、没有任何机制惩罚它们学到相似的东西，梯度也不会主动要求它们分散。真实模型里既有明确分工的头，也有大批彼此高度相似的头，还有退化成\emph{attention sink}的头——所有查询都把质量压在首个 token 上，等价于什么也不读。

所以本篇开头那个消融现象的正确解释不是``八个头各司其职''，而是弱一些也牢靠得多的一句：八条独立的低秩匹配通道，比一条高秩通道更容易被优化找到有用的解。表达力上界是数学，分工是不是真的发生，是经验问题。

顺带一提，图中两个头虽然读取规则迥异，却共享一个盲点：把输入序列整体重排，两张热度图会跟着一起重排，没有任何一个头会\emph{察觉}顺序变了。这个盲点是结构性的，下一篇会把它写成一条命题，并说明为什么位置信息必须从外面补进来。
